% !PRINTED 603
\weberpage{603}
Daraus ergiebt sich eine Basis von $\mathfrak{o}$:

a) \quad $1, \sqrt{d},$ \hfill $d\equiv 2$ oder $\equiv 3\ (\mathrm{mod}\ 4),$

b) \quad $1, \dfrac{1+\sqrt{d}}{2},$ \hfill $d\equiv 1\ (\mathrm{mod}\ 4).$

Wenn wir also

\weq{6}{\begin{aligned}
\text{a)}\quad \Theta &= \sqrt{d}, & d &\equiv 2, 3\ (\mathrm{mod}\ 4)\\
\text{b)}\quad \Theta &= \frac{1+\sqrt{d}}{2}, & d &\equiv 1\ (\mathrm{mod}\ 4)
\end{aligned}}
setzen, so ist $1, \Theta$ eine Basis von $\mathfrak{o}$.

F\"ur die Grundzahl erhalten wir demnach, wenn $\Theta'$ mit $\Theta$ conjugirt ist,
\[ \Delta = \left|\begin{array}{cc} 1, & \Theta \\ 1, & \Theta' \end{array}\right|^2 = (\Theta'-\Theta)^2, \]
also

\weq{7}{\begin{aligned}
\text{a)}\quad \Delta &= 4d, & d &\equiv 2, 3\ (\mathrm{mod}\ 4)\\
\text{b)}\quad \Delta &= d, & d &\equiv 1\ (\mathrm{mod}\ 4),
\end{aligned}}
\weq{8}{\Theta - \Theta' = \sqrt{\Delta}.}

Die Grundzahl ist also $\equiv 0$ oder $\equiv 1\ (\mathrm{mod}\ 4)$, und muss, wenn sie durch 4 theilbar ist, $\equiv 8, 12\ (\mathrm{mod}\ 16)$ sein\footnote{Eine ganze rationale Zahl $\Delta$, die nach dem Modul 4 mit 0 oder mit 1 congruent ist, hat \gesperrt{Kronecker} eine Zahl von Discriminantenform genannt. Hat $\Delta$ die Eigenschaft, dass sich kein quadratischer Factor so daraus absondern l\"asst, dass eine Zahl von Discriminantenform \"ubrig bleibt, so heisst $\Delta$ eine \gesperrt{Stammdiscriminante}. Die Grundzahl eines quadratischen K\"orpers ist also immer eine \gesperrt{Stammdiscriminante}, und umgekehrt ist auch jede Stammdiscriminante Grundzahl eines quadratischen K\"orpers. Vergl. \gesperrt{H. Weber}, Zahlentheoretische Untersuchungen aus dem Gebiete der elliptischen Functionen. I, II, III. Nachrichten der Gesellschaft der Wissenschaften zu G\"ottingen. 1893.}. Sie kann daher (in Uebereinstimmung mit dem Satze von \gesperrt{Minkowski}) weder $=\pm 1$ noch $=\pm 2$ sein. Die Grundzahl vom absolut kleinsten Werthe ist $\Delta = -3.$

\begin{center}
\S.\,165.\\
\gesperrt{Die Primideale in $\Omega$.}
\end{center}

Es bedeute nun $\varphi$ ein beliebiges ganzes Functional in $\Omega$. Wir wollen nach \S.\,146 eine Basis $\alpha_1, \alpha_2$ f\"ur $\varphi$, und damit die
%CONT>>
